% ---------------------------------------------------------
\section{Constant–Extraction Methodology (Optional)}
\label{app:C}

This appendix documents how numerical constants and slopes were extracted from the
experiments in Part~3.  All scripts and CSVs (\texttt{bias\_summary.csv},
\texttt{bias\_blocks.csv}, \texttt{variance\_table.csv}) are included in the repo and
reproduce Figures~\ref{fig:variance1}–\ref{fig:bias2} and the tables cited there.

\subsection*{Common settings}
Unless otherwise stated, blocks have
$\Delta=c_\Delta/\log T$ with $c_\Delta=2\pi$ (variance) or
adaptively chosen $c_\Delta\in[2\pi,60\pi]$ (bias),
and the penalty is $\lambda=c_\lambda(\log T)^{-2}$ with $c_\lambda\in\{1/4,1,4\}$.
Integrals on each block use the trapezoidal rule with at least~$12$ samples per block
(variance) or~$8$ (bias).
All runs were performed in \texttt{R 4.5.1} (double precision).

\paragraph{Data windows.}
Variance tests use Odlyzko ordinates $\gamma_k$ (on-line zeros, $\beta_k=\tfrac12$)
up to $T_{\max}=5\times10^4$.  In all experiments, only zeros with
$|\gamma_k|\le 200$ are retained (approximately $N(200)\approx 80$ ordinates);
this threshold suffices to resolve the harmonic field on the tested interval,
and the zero-count cap of $6{,}000$ is never binding.
Bias tests augment the on-line set by a synthetic
off-line zero at $\beta=\sigma+\eta$, with $\eta\in\{10^{-3},5\!\times\!10^{-4},10^{-4}\}$.

\medskip
\noindent\textbf{1) Derivative constant $c_{\mathrm{der}}$.}
We use the affine Poincaré inequality on a block $I$ of length $\Delta$:
\[
\int_I |r'(t)|^2\,dt \;\ge\; \frac{12}{\Delta^2}\int_I |r(t)|^2\,dt.
\]
Thus $c_{\mathrm{der}}=\tfrac{1}{12}$ in the normalization of Part~2
(see Appendix~\ref{app:B}).  Numerically, we verify the bound by
projection onto the affine subspace on a block of length~$\Delta$
and computing the ratio
$\|r'\|_{L^2(I)}^2 / \|r\|_{L^2(I)}^2$ over $10^4$ randomly phased test
signals $\sum_{m=1}^M a_m\cos(\omega_m t+\phi_m)$ with $M\in[1,8]$,
$\omega_m\Delta\in[0.1,10]$.  Every ratio exceeds the lower bound
$12/\Delta^2$; the empirical minimum lies near the sharp constant
$4\pi^2/\Delta^2\approx 39.5/\Delta^2$ (see Appendix~\ref{app:B}),
confirming that $c_{\mathrm{der}}=\tfrac1{12}$ is a valid but conservative
choice.

\medskip
\noindent\textbf{2) Variance constants $C_0,\,C_0'$.}
We bound
\(
\sum_j \int_{I_j} |H'_\sigma(t)|^2 dt
\le C_0\, T\log T\log\log T
\)
as in Appendix~\ref{app:A}.  From \texttt{variance\_table.csv}
(columns $F_\lambda$, $T\log T\log\log T$) we regress
$Y= \sum_j \int_{I_j} |H'_\sigma|^2$ on
$X=T\log T\log\log T$ \emph{through the origin} to obtain
$\widehat{C}_0$.  Reported intervals are nonparametric
percentile bootstrap ($B=2000$ resamples over the $T$-grid).
With only six $T$-values the bootstrap distribution is wide and
right-skewed: $\widehat{C}_0\approx 5.6\times10^{-4}$ with
$95\%$ CI $[6.7\times10^{-5},\,9.1\times10^{-3}]$,
reflecting the non-monotone behavior of~$F_\lambda$
across the tested range.  The affine residual bound
(Equation~\eqref{eq:aff-proj}) yields $C_0'=c_{\mathrm{aff}}^{-1}C_0$
with $c_{\mathrm{aff}}=\tfrac14$.

\medskip
\noindent\textbf{3) Slope calibration (bias regime).}
For each $(\eta,T)$ pair we compute the blockwise SPTB functional for
both the baseline (on-line) and biased (off-line zero injected) series.
We form the cumulative difference
$\Delta F_{\le j}=\sum_{k\le j}(F_{\lambda,k}^{\mathrm{bias}}-F_{\lambda,k}^{\mathrm{base}})$
and fit
\[
\log\bigl(\Delta F_{\le j}\bigr) \;=\; s\,t_j + b + \varepsilon
\]
over the last third of blocks (tail window)
using Huber regression (\texttt{MASS::rlm}, tuning $k=1.345$)
to reduce leverage from early-$t$ transients.
The differencing removes the $O(T\log T\log\log T)$ polynomial background,
so the extracted slope $\hat{s}$ estimates~$2\eta$.
Relative errors against the theoretical $2\eta$ are recorded in the
bias-regime tables of Part~3.
